% GENERATED FILE. Edit canonical lessons, then run npm run generate:books.
% canonical-source-hash: fnv1a64:1acd1d043576617e
% canonical-lessons: FR-C281-reportage, FR-C281-envoye-special, FR-C281-correspondant, FR-C281-depeche, FR-C281-agence-de-presse

\chapter{Le reportage}
\label{ch:fr-le-reportage}
\hlchaptermodality{281}

\begin{chapteropening}
\textbf{By the end of this chapter:} \emph{I can talk about reporting from the field: a report, a special correspondent, a correspondent, a wire story and a news agency.}

Complete the last lesson of chapter 281: I can talk about reporting from the field: a report, a special correspondent, a correspondent, a wire story and a news agency.
\end{chapteropening}

\section[le reportage]{le reportage — the report, the feature (on TV, radio or in a paper)}

\label{lesson:FR-C281-reportage}



\begin{quote}
Before the new one: say the French for the listener, then the French for the microphone.
\end{quote}

\subsection*{You'll want to know: le reportage}
\textbf{le reportage} — \textquotedblleft{}the report, the feature (on TV, radio or in a paper)\textquotedblright{}.

\begin{quote}
J'ai vu un reportage sur la vie à la montagne.
\end{quote}

\begin{quote}
\emph{I watched a feature about life in the mountains.}
\end{quote}

\emph{Faire un reportage} — to report on something, on location.

\begin{cousinweb}
\emph{reportage} is built with the suffix \emph{-age} on the verb meaning \textquotedblleft{}to carry back, to report\textquotedblright{} (EtymDB), which comes from Latin \emph{reportāre} \textquotedblleft{}to carry back\textquotedblright{}, built on \emph{portāre} \textquotedblleft{}to carry\textquotedblright{} (Lewis \& Short).
\end{cousinweb}

\begin{grammarlens}[title={What you've built}]
One more word for telling, explaining and describing at B1.
\end{grammarlens}

\subsection*{Your turn}
\begin{itemize}
\raggedright
  \item \emph{Say it:} \emph{le reportage}
  \item \emph{Say it:} \emph{le reportage}, once more
  \item \emph{Say it:} a sentence of your own with \emph{le reportage}
  \item \emph{Recall:} say the French for the listener, then the French for the microphone, then say \emph{le reportage} again
\end{itemize}

\begin{tcolorbox}[breakable,colback=teal!4,colframe=teal!35!black,title={Before you move on}]
What does le reportage mean? (\textquotedblleft{}The report, the feature (on TV, radio or in a paper)\textquotedblright{}.) Say it once more.
\end{tcolorbox}
\section[l'envoyé spécial]{l'envoyé spécial — the special correspondent}

\label{lesson:FR-C281-envoye-special}



\begin{quote}
Before the new one: say the French for the microphone, then the French for the report.
\end{quote}

\subsection*{You'll want to know: l'envoyé spécial}
\textbf{l'envoyé spécial} — \textquotedblleft{}the special correspondent\textquotedblright{}.

\begin{quote}
Notre envoyé spécial est arrivé sur place ce matin.
\end{quote}

\begin{quote}
\emph{Our special correspondent arrived on the spot this morning.}
\end{quote}

\emph{Envoyé} is the past participle of \emph{envoyer}: someone sent.

\begin{cousinweb}
\emph{envoyé} comes from \emph{envoyer}, inherited from Late Latin \emph{inviāre} \textquotedblleft{}to put on the way\textquotedblright{}, built on \emph{via} \textquotedblleft{}way\textquotedblright{} (Lewis \& Short); \emph{spécial} is a learned borrowing from Latin \emph{speciālis} (EtymDB).
\end{cousinweb}

\begin{grammarlens}[title={What you've built}]
One more word for telling, explaining and describing at B1.
\end{grammarlens}

\subsection*{Your turn}
\begin{itemize}
\raggedright
  \item \emph{Say it:} \emph{l'envoyé spécial}
  \item \emph{Say it:} \emph{l'envoyé spécial}, once more
  \item \emph{Say it:} a sentence of your own with \emph{l'envoyé spécial}
  \item \emph{Recall:} say the French for the microphone, then the French for the report, then say \emph{l'envoyé spécial} again
\end{itemize}

\begin{tcolorbox}[breakable,colback=teal!4,colframe=teal!35!black,title={Before you move on}]
What does l'envoyé spécial mean? (\textquotedblleft{}The special correspondent\textquotedblright{}.) Say it once more.
\end{tcolorbox}
\section[le correspondant]{le correspondant — the correspondent; a pen friend}

\label{lesson:FR-C281-correspondant}



\begin{quote}
Before the new one: say the French for the report, then the French for the special correspondent.
\end{quote}

\subsection*{You'll want to know: le correspondant}
\textbf{le correspondant} — \textquotedblleft{}the correspondent; a pen friend\textquotedblright{}.

\begin{quote}
Le correspondant du journal vit à l'étranger depuis dix ans.
\end{quote}

\begin{quote}
\emph{The paper's correspondent has lived abroad for ten years.}
\end{quote}

At school, \emph{un correspondant} is an exchange partner or a pen friend.

\begin{cousinweb}
\emph{correspondant} is built on \emph{correspondre}, a learned borrowing from Medieval Latin \emph{correspondēre} \textquotedblleft{}to answer together\textquotedblright{}, from \emph{con-} and \emph{respondēre} \textquotedblleft{}to answer\textquotedblright{} (Lewis \& Short).
\end{cousinweb}

\begin{grammarlens}[title={What you've built}]
One more word for telling, explaining and describing at B1.
\end{grammarlens}

\subsection*{Your turn}
\begin{itemize}
\raggedright
  \item \emph{Say it:} \emph{le correspondant}
  \item \emph{Say it:} \emph{le correspondant}, once more
  \item \emph{Say it:} a sentence of your own with \emph{le correspondant}
  \item \emph{Recall:} say the French for the report, then the French for the special correspondent, then say \emph{le correspondant} again
\end{itemize}

\begin{tcolorbox}[breakable,colback=teal!4,colframe=teal!35!black,title={Before you move on}]
What does le correspondant mean? (\textquotedblleft{}The correspondent; a pen friend\textquotedblright{}.) Say it once more.
\end{tcolorbox}
\section[la dépêche]{la dépêche — the dispatch, the wire story}

\label{lesson:FR-C281-depeche}



\begin{quote}
Before the new one: say the French for the special correspondent, then the French for the correspondent.
\end{quote}

\subsection*{You'll want to know: la dépêche}
\textbf{la dépêche} — \textquotedblleft{}the dispatch, the wire story\textquotedblright{}.

\begin{quote}
Une dépêche est arrivée à midi.
\end{quote}

\begin{quote}
\emph{A wire story came in at noon.}
\end{quote}

A short news item sent by an agency to the papers.

\begin{cousinweb}
\emph{dépêche} is built on \emph{dépêcher} \textquotedblleft{}to send off quickly\textquotedblright{}, which comes from Old French \emph{despeechier} \textquotedblleft{}to set free\textquotedblright{}, formed as the opposite of \emph{empêcher} \textquotedblleft{}to hinder\textquotedblright{}; the root is usually traced to Latin \emph{pedica} \textquotedblleft{}a fetter\textquotedblright{} (Lewis \& Short).
\end{cousinweb}

\begin{grammarlens}[title={What you've built}]
One more word for telling, explaining and describing at B1.
\end{grammarlens}

\subsection*{Your turn}
\begin{itemize}
\raggedright
  \item \emph{Say it:} \emph{la dépêche}
  \item \emph{Say it:} \emph{la dépêche}, once more
  \item \emph{Say it:} a sentence of your own with \emph{la dépêche}
  \item \emph{Recall:} say the French for the special correspondent, then the French for the correspondent, then say \emph{la dépêche} again
\end{itemize}

\begin{tcolorbox}[breakable,colback=teal!4,colframe=teal!35!black,title={Before you move on}]
What does la dépêche mean? (\textquotedblleft{}The dispatch, the wire story\textquotedblright{}.) Say it once more.
\end{tcolorbox}
\section[l'agence de presse]{l'agence de presse — the news agency}

\label{lesson:FR-C281-agence-de-presse}



\begin{quote}
Before the new one: say the French for the correspondent, then the French for the dispatch.
\end{quote}

\subsection*{You'll want to know: l'agence de presse}
\textbf{l'agence de presse} — \textquotedblleft{}the news agency\textquotedblright{}.

\begin{quote}
L'agence de presse envoie des dépêches toute la journée.
\end{quote}

\begin{quote}
\emph{The news agency sends out wire stories all day long.}
\end{quote}

Other agencies: \emph{une agence de voyages}, \emph{une agence immobilière}.

\begin{cousinweb}
\emph{agence} is usually traced to Italian \emph{agenzia} (EtymDB lists no parent), which goes back to Latin \emph{agere} \textquotedblleft{}to do, to act\textquotedblright{} (Lewis \& Short); \emph{presse} is built on \emph{presser}, from Latin \emph{premere} \textquotedblleft{}to press\textquotedblright{}.
\end{cousinweb}

\begin{grammarlens}[title={What you've built}]
One more word for telling, explaining and describing at B1.
\end{grammarlens}

\subsection*{Your turn}
\begin{itemize}
\raggedright
  \item \emph{Say it:} \emph{l'agence de presse}
  \item \emph{Say it:} \emph{l'agence de presse}, once more
  \item \emph{Say it:} a sentence of your own with \emph{l'agence de presse}
  \item \emph{Recall:} say the French for the correspondent, then the French for the dispatch, then say \emph{l'agence de presse} again
\end{itemize}

\begin{tcolorbox}[breakable,colback=teal!4,colframe=teal!35!black,title={Before you move on}]
What does l'agence de presse mean? (\textquotedblleft{}The news agency\textquotedblright{}.) Say it once more.
\end{tcolorbox}
